\documentclass[10pt,a4paper]{amsart}
\usepackage[margin=19mm]{geometry}
\usepackage{amsmath,amssymb,mathtools}
\usepackage[scr=rsfs,cal=euler]{mathalfa}
\usepackage{xfrac}
\usepackage[unicode,hidelinks,bookmarks=false]{hyperref}
\newcommand{\ie}{i.e.\ }
\newcommand{\cf}{cf.\ }
\newcommand{\resp}{respectively\ }
\newcommand{\Subsection}{\subsection}
\providecommand{\nobreakdash}{\nobreak\hbox{-}}
\setlength{\emergencystretch}{2em}
\multlinegap0pt
\usepackage{amssymb}
\usepackage{xcolor}
\newtheorem{lemma}{Lemma}
\newtheorem{theorem}{Theorem}
\newcommand{\Ha}{\mathcal{H}}
\newcommand{\I}{I_N}
\newcommand{\N}{\mathbb{N}}
\DeclareMathOperator{\dimH}{dim_H}
\title{Sharp Hausdorff Dimension Bounds for Sets with Bounded and Growing Digits in $N$-expansions}
\author{Andreea Catalina Chitu, Gabriela Ileana Sebe, Dan Lascu}
\date{Source: arXiv:2603.27635v1 (CC BY 4.0)}
\begin{document}
\maketitle

\noindent\emph{Source and licence.} Sharp Hausdorff Dimension Bounds for Sets with Bounded and Growing Digits in $N$-expansions. Version arXiv:2603.27635v1 (CC BY 4.0), distributed by arXiv under the Creative Commons Attribution 4.0 International licence, http://creativecommons.org/licenses/by/4.0/, as recorded in the arXiv licence metadata for this exact version and shown on the abstract page https://arxiv.org/abs/2603.27635v1. Full licence notice: \texttt{licenses/arxiv-2603.27635-CC-BY-4.0.txt}.\par\smallskip

\noindent\textit{Editorial scope.} Counted unit: the article's theorem upperbound:Good2N (source0.tex lines 769--776). The article's complete original proof of it is delivered with it (source0.tex lines 777--894, one \texttt{proof} environment of 118 lines). No mathematics is written, repaired or re-derived: the statement and the proof are the article's own lines, in the article's own notation, and no retained line is substituted or re-flowed.  Retained uncounted as the source-local context the proof needs: lines 199--206; lines 296--298.  Nothing was omitted from any retained span, so the package carries no omission record.  No retained line contains a citation, so the wrapper carries no bibliography and no bibliography entry.  No retained line contains a figure environment, an image-inclusion command or drawing code, so no image is delivered and the package declares no asset dependency.  Editorial formatting only, in addition: an AMS \texttt{amsart} preamble that restates the article's own declarations and macros used by the retained lines, namely the environments \texttt{lemma}, \texttt{theorem} on separate counters, together with the standard packages those lines need.  The retained environments print in this document as Lemma 1, Theorem 1 in order of appearance, whereas the article prints the same items with the numbering of its own full document.  The complete native source of the article as distributed in the arXiv e-print for this exact version is bundled unchanged as \texttt{sources/197366/source0.tex}.
\begin{lemma} \label{lemma:q-growth}
For all $n \in \N_+$ and sequences $\varepsilon_1, \ldots, \varepsilon_n \geq N$, the denominators satisfy:
\begin{enumerate}
\item[(i)] ${q_{n}(\varepsilon_1, \ldots, \varepsilon_n)} \ge N {q_{n-1}(\varepsilon_1, \ldots, \varepsilon_{n-1})}$;
\item[(ii)] $q_n(\varepsilon_1, \ldots, \varepsilon_n) \geq N^n$;
\item[(iii)] $q_n(\varepsilon_1, \ldots, \varepsilon_n) \geq (2N)^{\frac{n-1}{2}}$.
\end{enumerate}
\end{lemma}

\begin{equation}
\frac{N^{n+1}}{(1+N)q_n^2} \leq |I_n(\varepsilon_1, \ldots, \varepsilon_n)| \leq \frac{N^n}{q_n^2}; \label{eq:interval-bounds}
\end{equation}

\begin{theorem}[Upper Bound for Growing Digits] \label{upperbound:Good2N}
Let $N \geq 1$ be an integer. For every real $\alpha$ with $\log(\alpha-1) > \exp(1+N)$, we have
\begin{equation}\label{eq:upper-ineq}
\dimH(F_{\alpha,N}) 
\;\le\; 
\frac12 + \frac{\log\log\!\bigl({\alpha}-1\bigr)}{2\log\!\bigl({\alpha}-1\bigr)}.
\end{equation}
\end{theorem}

\begin{proof} 
Let $s > \tfrac12$ and for each $n \geq 1$ consider the family
\[
\mathcal{C}_n := \{ \I(\varepsilon_1, \dots, \varepsilon_n) : \varepsilon_i \geq \alpha,\; 1 \leq i \leq n \}.
\]
By Lemma~\ref{lemma:q-growth} and the upper bound in \eqref{eq:interval-bounds}, 
the diameters of intervals in $\mathcal{C}_n$ tend to $0$ as $n \to \infty$. Hence $\mathcal{C}_n$ is a cover of $F_{\alpha,N}$.

For a fixed $I_n = \I(\varepsilon_1,\dots,\varepsilon_n) \in \mathcal{C}_n$, its children are 
$I_{n+1}^{(k)} = \I(\varepsilon_1,\dots,\varepsilon_n,k)$ with $k \geq \alpha$. 
Using $q_{n+1} \geq k q_n$ and the upper bound in \eqref{eq:interval-bounds},
\begin{equation}
\sum_{k=\alpha}^{\infty} |I_{n+1}^{(k)}|^s 
   \leq \sum_{k=\alpha}^{\infty} \left( \frac{N^{\,n+1}}{k^2 q_n^{2}} \right)^s
   = \frac{N^{(n+1)s}}{q_n^{2s}} \sum_{k=\alpha}^{\infty} k^{-2s}.
\label{eq:sum-children}
\end{equation}

On the other hand, the lower bound in \eqref{eq:interval-bounds} gives
\begin{equation}\label{eq:parent-bound}
|I_n|^s \geq \left( \frac{N^{\,n+1}}{(1+N) q_n^{2}} \right)^s 
        = \frac{N^{(n+1)s}}{(1+N)^s q_n^{2s}}.
\end{equation}

From \eqref{eq:sum-children} and \eqref{eq:parent-bound} we obtain the key inequality:
\begin{equation}\label{eq:children-parent}
\sum_{k=\alpha}^{\infty} |I_{n+1}^{(k)}|^s 
\; \leq \; (1+N)^s \left( \sum_{k=\alpha}^{\infty} k^{-2s} \right) |I_n|^s .
\end{equation}

This inequality shows that the $s$-dimensional sum over all children of $I_n$ is bounded by 
$(1+N)^s \left( \sum_{k=\alpha}^{\infty} k^{-2s} \right)$ times the $s$-dimensional size of the parent interval.

If we have
\begin{equation}\label{eq:key-condition}
\sum_{k=\alpha}^{\infty} k^{-2s} \leq (1+N)^{-s},
\end{equation}
then from \eqref{eq:children-parent} we obtain for every $I_n \in \mathcal{C}_n$,
\begin{equation}\label{eq:per-parent}
\sum_{k=\alpha}^{\infty} |I_{n+1}^{(k)}|^s \leq |I_n|^s .
\end{equation}

We now sum inequality \eqref{eq:per-parent} over all parent intervals $I_n \in \mathcal{C}_n$. 
On the left-hand side, we obtain a double sum:
\[
\sum_{I_n \in \mathcal{C}_n} \sum_{k=\alpha}^{\infty} |I_{n+1}^{(k)}|^s .
\]

Observe that each interval $I_{n+1} \in \mathcal{C}_{n+1}$ corresponds to a unique pair 
$(I_n, k)$ where $I_n \in \mathcal{C}_n$ is the parent interval and $k \geq \alpha$ is the 
$(n+1)$-th digit. This is because each fundamental interval of order $n+1$ has a unique 
representation as $I_N(\varepsilon_1, \dots, \varepsilon_n, k)$. 

Therefore, the double sum over all parents $I_n$ and all digits $k \geq \alpha$ is exactly 
the sum over all children intervals $I_{n+1} \in \mathcal{C}_{n+1}$:
\[
\sum_{I_n \in \mathcal{C}_n} \sum_{k=\alpha}^{\infty} |I_{n+1}^{(k)}|^s 
= \sum_{I_{n+1} \in \mathcal{C}_{n+1}} |I_{n+1}|^s .
\]

Summing the right-hand side of \eqref{eq:per-parent} over all $I_n \in \mathcal{C}_n$ gives 
simply $\sum_{I_n \in \mathcal{C}_n} |I_n|^s$.

Consequently, from \eqref{eq:per-parent} we obtain
\begin{equation}\label{eq:level-sum}
\sum_{I_{n+1} \in \mathcal{C}_{n+1}} |I_{n+1}|^s \leq \sum_{I_n \in \mathcal{C}_n} |I_n|^s .
\end{equation}

Define $a_n := \sum_{I \in \mathcal{C}_n} |I|^s$. Inequality \eqref{eq:level-sum} tells us that 
the sequence $(a_n)_{n\geq 1}$ is non‑increasing. In particular, $a_n \leq a_1$ for every $n$.

For any $\delta > 0$, choose $n$ large enough so that $\max\{|I| : I \in \mathcal{C}_n\} < \delta$
(this is possible because the diameters of the intervals in $\mathcal{C}_n$ tend to $0$ as $n$ increases). 
Then $\mathcal{C}_n$ is a $\delta$‑cover of $F_{\alpha,N}$ and
\[
\sum_{I \in \mathcal{C}_n} |I|^s = a_n \leq a_1 .
\]
Hence $\Ha^s_\delta(F_{\alpha,N}) \leq a_1$. Since $\delta$ was arbitrary, we obtain 
$\Ha^s(F_{\alpha,N}) \leq a_1 < \infty$, and consequently $\dimH(F_{\alpha,N}) \leq s$.

To obtain an explicit $s$ satisfying \eqref{eq:key-condition}, we approximate the sum by an integral.
For $s > \tfrac12$, the function $x \mapsto x^{-2s}$ is decreasing on $(0,\infty)$, hence
\[
\sum_{k=\alpha}^{\infty} k^{-2s} \leq \int_{\alpha-1}^{\infty} x^{-2s} \, dx.
\]

Evaluating the integral:
\[
\int_{\alpha-1}^{\infty} x^{-2s} \, dx = \left[ \frac{x^{-(2s-1)}}{-(2s-1)} \right]_{\alpha-1}^{\infty} 
= \frac{(\alpha-1)^{-(2s-1)}}{2s-1}.
\]
Since  $(1+N)^{-1} \le (1+N)^{-s}$ for $s \in (\frac{1}{2}, 1]$ it follows that a sufficient condition for \eqref{eq:key-condition} is:
\[
\frac{(\alpha-1)^{-(2s-1)}}{2s-1} \leq (1+N)^{-1},
\]
which is equivalent to 
$(2s-1)(\alpha-1)^{2s-1} \ge 1+N$. 

Therefore, $\dimH F_{N,\alpha} \le s$ where
\begin{equation} \label{5.6}
(2s-1)(\alpha-1)^{2s-1} = 1+N.
\end{equation}
This equation gives a unique value for $s \in (\frac{1}{2}, 1]$, if $\alpha \ge N+2$. 

For $s \ge \frac{1}{2}$, the left-hand side is a strictly increasing function of $2s-1$ which vanishes when $s=\frac{1}{2}$ and exceeds the right-hand side when $s=1$.

Putting $2s-1=x$ and $\alpha -1 =a$, equation \eqref{5.6} becomes 
\begin{equation} \label{5.7}
    x a^x = 1+N.
\end{equation}
When $x = \dfrac{\log\log a}{\log a}$, the left-hand side of \eqref{5.7} reduces to $\log \log a$ and this is greater than $1+N$ if $a > \exp(\exp(1+N))$. 
Hence the solution of \eqref{5.7} is less than $\dfrac{\log\log a}{\log a}$ if $a > \exp(\exp(1+N))$. 
Therefore, 
\[
\dimH F_{N,\alpha} \le \frac{1}{2} + \dfrac{\log\log (\alpha-1)}{2}\log (\alpha-1)
\]
if $\log(\alpha-1) > \exp(1+N)$. 
\end{proof}

\end{document}
